
\documentclass{article}
\usepackage{colortbl}
\usepackage{makecell}
\usepackage{multirow}
\usepackage{supertabular}

\begin{document}

\newcounter{utterance}

\centering \large Interaction Transcript for game `chronicle', experiment `battles\_easy', episode 0 with qwen3.5{-}4b.
\vspace{24pt}

{ \footnotesize  \setcounter{utterance}{1}
\setlength{\tabcolsep}{0pt}
\begin{supertabular}{c@{$\;$}|p{.15\linewidth}@{}p{.15\linewidth}p{.15\linewidth}p{.15\linewidth}p{.15\linewidth}p{.15\linewidth}}
    \# & $\;$A & \multicolumn{4}{c}{Game Master} & $\;\:$B\\
    \hline

    \theutterance \stepcounter{utterance}  
    & & \multicolumn{4}{p{0.6\linewidth}}{
        \cellcolor[rgb]{0.9,0.9,0.9}{
            \makecell[{{p{\linewidth}}}]{
                \texttt{\tiny{[P1$\langle$GM]}}
                \texttt{You are playing a historical deduction game as the NARRATOR.} \\
\\ 
\texttt{You have been given the name of a historical event, discovery, or milestone. Your goal is to help the DETECTIVE figure out which event you are describing — but you must describe it indirectly, without ever naming it.} \\
\\ 
\texttt{Rules:} \\
\texttt{1. Each round, write exactly ONE sentence describing the event. Add NEW information each round — do not repeat what you wrote before.} \\
\texttt{2. You must NEVER use any of the following forbidden keywords (these would make the event too easy to identify): iran, persian, russian, russo, safavid} \\
\texttt{3. You must NEVER state the event name directly or any close variant of it.} \\
\texttt{4. You may freely describe causes, consequences, the broader context, the sequence of events, and the significance — but only using general, non{-}identifying language.} \\
\\ 
\texttt{Format your response EXACTLY as follows (required):} \\
\texttt{PARAGRAPH: <exactly one sentence describing the event without forbidden keywords>} \\
\\ 
\texttt{The event you must describe (without naming it) is: Russo{-}Persian War (1722–1723)} \\
            }
        }
    }
    & \\ \\

    \theutterance \stepcounter{utterance}  
    & \multicolumn{4}{p{0.6\linewidth}}{
        \cellcolor[rgb]{0.9,0.9,0.9}{
            \makecell[{{p{\linewidth}}}]{
                \texttt{\tiny{[P1$\rangle$GM]}}
                \texttt{PARAGRAPH: This conflict marked the first major military engagement between the expanding empire of the north and the newly established dynasty in the east, resulting in a decisive victory for the latter and the consolidation of their control over the region.} \\
            }
        }
    }
    & & \\ \\

    \theutterance \stepcounter{utterance}  
    & & \multicolumn{4}{p{0.6\linewidth}}{
        \cellcolor[rgb]{0.9,0.9,0.9}{
            \makecell[{{p{\linewidth}}}]{
                \texttt{\tiny{[GM$\rangle$P2]}}
                \texttt{You are playing a historical deduction game as the DETECTIVE.} \\
\\ 
\texttt{A Narrator will describe a historical event, discovery, or milestone to you over several rounds. Each round you receive a new sentence of description that builds on what came before. Your job is to figure out which specific event is being described.} \\
\\ 
\texttt{After EACH new sentence, respond briefly. Keep your entire response under 100 words. Use this exact format (both sections are required):} \\
\\ 
\texttt{ANALYSIS: <2–3 short sentences: what does the growing description suggest? What type of event is this? What are the key features you are noticing?>} \\
\\ 
\texttt{GUESS: <your best guess for the name of the event being described — be specific, e.g. "The French Revolution" not just "a revolution">} \\
\\ 
\texttt{Rules:} \\
\texttt{{-} You may refine your GUESS each round as more information arrives.} \\
\texttt{{-} Your GUESS should name the specific event, not just a broad category.} \\
\texttt{{-} You have up to 10 rounds total.} \\
\texttt{{-} Be concise — quality of reasoning matters more than length.} \\
\\ 
\texttt{The first sentence from the Narrator is coming now.} \\
\\ 
\\ 
\texttt{{[}Sentence 1{]}} \\
\texttt{This conflict marked the first major military engagement between the expanding empire of the north and the newly established dynasty in the east, resulting in a decisive victory for the latter and the consolidation of their control over the region.} \\
            }
        }
    }
    & \\ \\

    \theutterance \stepcounter{utterance}  
    & & & \multicolumn{4}{p{0.6\linewidth}}{
        \cellcolor[rgb]{0.9,0.9,0.9}{
            \makecell[{{p{\linewidth}}}]{
                \texttt{\tiny{[GM$\langle$P2]}}
                \texttt{ANALYSIS: The description suggests a military conflict between a northern empire and an eastern dynasty, resulting in a decisive victory that consolidated eastern control. This points to a specific historical war involving territorial expansion and political unification.} \\
\\ 
\texttt{GUESS: The Battle of Red Cliffs} \\
            }
        }
    }
     \\ \\

    \theutterance \stepcounter{utterance}  
    & & \multicolumn{4}{p{0.6\linewidth}}{
        \cellcolor[rgb]{0.9,0.9,0.9}{
            \makecell[{{p{\linewidth}}}]{
                \texttt{\tiny{[P1$\langle$GM]}}
                \texttt{Round 1 complete. The Detective's current guess: "The Battle of Red Cliffs". Please write a new paragraph with fresh information.} \\
            }
        }
    }
    & \\ \\

    \theutterance \stepcounter{utterance}  
    & \multicolumn{4}{p{0.6\linewidth}}{
        \cellcolor[rgb]{0.9,0.9,0.9}{
            \makecell[{{p{\linewidth}}}]{
                \texttt{\tiny{[P1$\rangle$GM]}}
                \texttt{PARAGRAPH: The war concluded with a treaty that formally recognized the independence of the eastern kingdom and established a long{-}standing border between the two great powers.} \\
            }
        }
    }
    & & \\ \\

    \theutterance \stepcounter{utterance}  
    & & \multicolumn{4}{p{0.6\linewidth}}{
        \cellcolor[rgb]{0.9,0.9,0.9}{
            \makecell[{{p{\linewidth}}}]{
                \texttt{\tiny{[GM$\rangle$P2]}}
                \texttt{{[}Sentence 2{]}} \\
\texttt{The war concluded with a treaty that formally recognized the independence of the eastern kingdom and established a long{-}standing border between the two great powers.} \\
            }
        }
    }
    & \\ \\

    \theutterance \stepcounter{utterance}  
    & & & \multicolumn{4}{p{0.6\linewidth}}{
        \cellcolor[rgb]{0.9,0.9,0.9}{
            \makecell[{{p{\linewidth}}}]{
                \texttt{\tiny{[GM$\langle$P2]}}
                \texttt{ANALYSIS: The mention of a treaty recognizing independence and establishing a border suggests a conflict that ended without total conquest, likely involving a balance of power. The previous guess of Red Cliffs fits the "eastern dynasty" (Wei) vs "northern empire" (Cao Cao) dynamic, though the treaty detail is specific.} \\
\\ 
\texttt{GUESS: The Battle of Red Cliffs} \\
            }
        }
    }
     \\ \\

    \theutterance \stepcounter{utterance}  
    & & \multicolumn{4}{p{0.6\linewidth}}{
        \cellcolor[rgb]{0.9,0.9,0.9}{
            \makecell[{{p{\linewidth}}}]{
                \texttt{\tiny{[P1$\langle$GM]}}
                \texttt{Round 2 complete. The Detective's current guess: "The Battle of Red Cliffs". Please write a new paragraph with fresh information.} \\
            }
        }
    }
    & \\ \\

    \theutterance \stepcounter{utterance}  
    & \multicolumn{4}{p{0.6\linewidth}}{
        \cellcolor[rgb]{0.9,0.9,0.9}{
            \makecell[{{p{\linewidth}}}]{
                \texttt{\tiny{[P1$\rangle$GM]}}
                \texttt{PARAGRAPH: This prolonged struggle was driven by the intense rivalry between two neighboring superpowers vying for dominance over the trade routes and territories stretching from the Caspian Sea to the Indian Ocean.} \\
            }
        }
    }
    & & \\ \\

    \theutterance \stepcounter{utterance}  
    & & \multicolumn{4}{p{0.6\linewidth}}{
        \cellcolor[rgb]{0.9,0.9,0.9}{
            \makecell[{{p{\linewidth}}}]{
                \texttt{\tiny{[GM$\rangle$P2]}}
                \texttt{{[}Sentence 3{]}} \\
\texttt{This prolonged struggle was driven by the intense rivalry between two neighboring superpowers vying for dominance over the trade routes and territories stretching from the Caspian Sea to the Indian Ocean.} \\
            }
        }
    }
    & \\ \\

    \theutterance \stepcounter{utterance}  
    & & & \multicolumn{4}{p{0.6\linewidth}}{
        \cellcolor[rgb]{0.9,0.9,0.9}{
            \makecell[{{p{\linewidth}}}]{
                \texttt{\tiny{[GM$\langle$P2]}}
                \texttt{ANALYSIS: The description shifts from a single battle to a prolonged struggle over trade routes between the Caspian and Indian Oceans. This geographic scope and the nature of the rivalry strongly suggest the conflict between the Roman Empire and the Sassanid Empire.} \\
\\ 
\texttt{GUESS: The Roman{-}Sassanid Wars} \\
            }
        }
    }
     \\ \\

    \theutterance \stepcounter{utterance}  
    & & \multicolumn{4}{p{0.6\linewidth}}{
        \cellcolor[rgb]{0.9,0.9,0.9}{
            \makecell[{{p{\linewidth}}}]{
                \texttt{\tiny{[P1$\langle$GM]}}
                \texttt{Round 3 complete. The Detective's current guess: "The Roman{-}Sassanid Wars". Please write a new paragraph with fresh information.} \\
            }
        }
    }
    & \\ \\

    \theutterance \stepcounter{utterance}  
    & \multicolumn{4}{p{0.6\linewidth}}{
        \cellcolor[rgb]{0.9,0.9,0.9}{
            \makecell[{{p{\linewidth}}}]{
                \texttt{\tiny{[P1$\rangle$GM]}}
                \texttt{PARAGRAPH: The outcome of this campaign significantly altered the geopolitical balance in the region, leading to a period of relative peace and the eventual rise of a new imperial dynasty in the west.} \\
            }
        }
    }
    & & \\ \\

    \theutterance \stepcounter{utterance}  
    & & \multicolumn{4}{p{0.6\linewidth}}{
        \cellcolor[rgb]{0.9,0.9,0.9}{
            \makecell[{{p{\linewidth}}}]{
                \texttt{\tiny{[GM$\rangle$P2]}}
                \texttt{{[}Sentence 4{]}} \\
\texttt{The outcome of this campaign significantly altered the geopolitical balance in the region, leading to a period of relative peace and the eventual rise of a new imperial dynasty in the west.} \\
            }
        }
    }
    & \\ \\

    \theutterance \stepcounter{utterance}  
    & & & \multicolumn{4}{p{0.6\linewidth}}{
        \cellcolor[rgb]{0.9,0.9,0.9}{
            \makecell[{{p{\linewidth}}}]{
                \texttt{\tiny{[GM$\langle$P2]}}
                \texttt{ANALYSIS: The shift from a prolonged war to a new western dynasty suggests a decisive turning point. The mention of a new dynasty in the west following a major campaign points toward the fall of Rome and the rise of the Franks.} \\
\\ 
\texttt{GUESS: The Battle of Tours} \\
            }
        }
    }
     \\ \\

    \theutterance \stepcounter{utterance}  
    & & \multicolumn{4}{p{0.6\linewidth}}{
        \cellcolor[rgb]{0.9,0.9,0.9}{
            \makecell[{{p{\linewidth}}}]{
                \texttt{\tiny{[P1$\langle$GM]}}
                \texttt{Round 4 complete. The Detective's current guess: "The Battle of Tours". Please write a new paragraph with fresh information.} \\
            }
        }
    }
    & \\ \\

    \theutterance \stepcounter{utterance}  
    & \multicolumn{4}{p{0.6\linewidth}}{
        \cellcolor[rgb]{0.9,0.9,0.9}{
            \makecell[{{p{\linewidth}}}]{
                \texttt{\tiny{[P1$\rangle$GM]}}
                \texttt{PARAGRAPH: The conflict ended with a decisive victory for the eastern forces, which secured their sovereignty against the western expansion and solidified their status as a major regional power.} \\
            }
        }
    }
    & & \\ \\

    \theutterance \stepcounter{utterance}  
    & & \multicolumn{4}{p{0.6\linewidth}}{
        \cellcolor[rgb]{0.9,0.9,0.9}{
            \makecell[{{p{\linewidth}}}]{
                \texttt{\tiny{[GM$\rangle$P2]}}
                \texttt{{[}Sentence 5{]}} \\
\texttt{The conflict ended with a decisive victory for the eastern forces, which secured their sovereignty against the western expansion and solidified their status as a major regional power.} \\
            }
        }
    }
    & \\ \\

    \theutterance \stepcounter{utterance}  
    & & & \multicolumn{4}{p{0.6\linewidth}}{
        \cellcolor[rgb]{0.9,0.9,0.9}{
            \makecell[{{p{\linewidth}}}]{
                \texttt{\tiny{[GM$\langle$P2]}}
                \texttt{ANALYSIS: The description now emphasizes a decisive eastern victory securing sovereignty against western expansion. This contradicts the previous guess of Tours (where the Franks won) and suggests a different conflict, possibly involving the Byzantine Empire or a specific campaign like the Battle of Manzikert.} \\
\\ 
\texttt{GUESS: The Battle of Manzikert} \\
            }
        }
    }
     \\ \\

    \theutterance \stepcounter{utterance}  
    & & \multicolumn{4}{p{0.6\linewidth}}{
        \cellcolor[rgb]{0.9,0.9,0.9}{
            \makecell[{{p{\linewidth}}}]{
                \texttt{\tiny{[P1$\langle$GM]}}
                \texttt{Round 5 complete. The Detective's current guess: "The Battle of Manzikert". Please write a new paragraph with fresh information.} \\
            }
        }
    }
    & \\ \\

    \theutterance \stepcounter{utterance}  
    & \multicolumn{4}{p{0.6\linewidth}}{
        \cellcolor[rgb]{0.9,0.9,0.9}{
            \makecell[{{p{\linewidth}}}]{
                \texttt{\tiny{[P1$\rangle$GM]}}
                \texttt{PARAGRAPH: This engagement was the culmination of decades of border skirmishes that ultimately forced the western empire to retreat from its eastern territories, leaving the eastern dynasty to expand its influence into the Caucasus and Central Asia.} \\
            }
        }
    }
    & & \\ \\

    \theutterance \stepcounter{utterance}  
    & & \multicolumn{4}{p{0.6\linewidth}}{
        \cellcolor[rgb]{0.9,0.9,0.9}{
            \makecell[{{p{\linewidth}}}]{
                \texttt{\tiny{[GM$\rangle$P2]}}
                \texttt{{[}Sentence 6{]}} \\
\texttt{This engagement was the culmination of decades of border skirmishes that ultimately forced the western empire to retreat from its eastern territories, leaving the eastern dynasty to expand its influence into the Caucasus and Central Asia.} \\
            }
        }
    }
    & \\ \\

    \theutterance \stepcounter{utterance}  
    & & & \multicolumn{4}{p{0.6\linewidth}}{
        \cellcolor[rgb]{0.9,0.9,0.9}{
            \makecell[{{p{\linewidth}}}]{
                \texttt{\tiny{[GM$\langle$P2]}}
                \texttt{ANALYSIS: The description confirms a decisive eastern victory forcing a western retreat, specifically mentioning expansion into the Caucasus and Central Asia. This aligns perfectly with the Battle of Manzikert (1071), where Seljuk Turks defeated the Byzantine Empire, opening the way for their expansion into Anatolia and beyond.} \\
\\ 
\texttt{GUESS: The Battle of Manzikert} \\
            }
        }
    }
     \\ \\

    \theutterance \stepcounter{utterance}  
    & & \multicolumn{4}{p{0.6\linewidth}}{
        \cellcolor[rgb]{0.9,0.9,0.9}{
            \makecell[{{p{\linewidth}}}]{
                \texttt{\tiny{[P1$\langle$GM]}}
                \texttt{Round 6 complete. The Detective's current guess: "The Battle of Manzikert". Please write a new paragraph with fresh information.} \\
            }
        }
    }
    & \\ \\

    \theutterance \stepcounter{utterance}  
    & \multicolumn{4}{p{0.6\linewidth}}{
        \cellcolor[rgb]{0.9,0.9,0.9}{
            \makecell[{{p{\linewidth}}}]{
                \texttt{\tiny{[P1$\rangle$GM]}}
                \texttt{PARAGRAPH: The aftermath of this war saw the eastern power successfully repel the northern invasions, allowing them to consolidate their hold over the Iranian plateau and secure their position as the dominant force in the region.} \\
            }
        }
    }
    & & \\ \\

    \theutterance \stepcounter{utterance}  
    & & \multicolumn{4}{p{0.6\linewidth}}{
        \cellcolor[rgb]{0.9,0.9,0.9}{
            \makecell[{{p{\linewidth}}}]{
                \texttt{\tiny{[GM$\rangle$P2]}}
                \texttt{{[}Sentence 7{]}} \\
\texttt{The aftermath of this war saw the eastern power successfully repel the northern invasions, allowing them to consolidate their hold over the Iranian plateau and secure their position as the dominant force in the region.} \\
            }
        }
    }
    & \\ \\

    \theutterance \stepcounter{utterance}  
    & & & \multicolumn{4}{p{0.6\linewidth}}{
        \cellcolor[rgb]{0.9,0.9,0.9}{
            \makecell[{{p{\linewidth}}}]{
                \texttt{\tiny{[GM$\langle$P2]}}
                \texttt{ANALYSIS: The description shifts focus from the Byzantine{-}Seljuk conflict to a different era involving the repelling of northern invasions and consolidation over the Iranian plateau. This suggests a conflict against the Mongols or a later Safavid victory, rather than Manzikert.} \\
\\ 
\texttt{GUESS: The Battle of Chaldiran} \\
            }
        }
    }
     \\ \\

    \theutterance \stepcounter{utterance}  
    & & \multicolumn{4}{p{0.6\linewidth}}{
        \cellcolor[rgb]{0.9,0.9,0.9}{
            \makecell[{{p{\linewidth}}}]{
                \texttt{\tiny{[P1$\langle$GM]}}
                \texttt{Round 7 complete. The Detective's current guess: "The Battle of Chaldiran". Please write a new paragraph with fresh information.} \\
            }
        }
    }
    & \\ \\

    \theutterance \stepcounter{utterance}  
    & \multicolumn{4}{p{0.6\linewidth}}{
        \cellcolor[rgb]{0.9,0.9,0.9}{
            \makecell[{{p{\linewidth}}}]{
                \texttt{\tiny{[P1$\rangle$GM]}}
                \texttt{PARAGRAPH: The final resolution of this struggle involved the signing of a peace accord that defined the boundaries of the two empires and ended a era of continuous warfare between them.} \\
            }
        }
    }
    & & \\ \\

    \theutterance \stepcounter{utterance}  
    & & \multicolumn{4}{p{0.6\linewidth}}{
        \cellcolor[rgb]{0.9,0.9,0.9}{
            \makecell[{{p{\linewidth}}}]{
                \texttt{\tiny{[GM$\rangle$P2]}}
                \texttt{{[}Sentence 8{]}} \\
\texttt{The final resolution of this struggle involved the signing of a peace accord that defined the boundaries of the two empires and ended a era of continuous warfare between them.} \\
            }
        }
    }
    & \\ \\

    \theutterance \stepcounter{utterance}  
    & & & \multicolumn{4}{p{0.6\linewidth}}{
        \cellcolor[rgb]{0.9,0.9,0.9}{
            \makecell[{{p{\linewidth}}}]{
                \texttt{\tiny{[GM$\langle$P2]}}
                \texttt{ANALYSIS: The description points to a formal peace treaty defining borders between two empires after a long era of war. This fits the conclusion of the Safavid{-}Ottoman conflicts, specifically the Treaty of Amasya or the Treaty of Zuhab, which ended decades of fighting and established the modern borders of the region.} \\
\\ 
\texttt{GUESS: The Treaty of Zuhab} \\
            }
        }
    }
     \\ \\

    \theutterance \stepcounter{utterance}  
    & & \multicolumn{4}{p{0.6\linewidth}}{
        \cellcolor[rgb]{0.9,0.9,0.9}{
            \makecell[{{p{\linewidth}}}]{
                \texttt{\tiny{[P1$\langle$GM]}}
                \texttt{Round 8 complete. The Detective's current guess: "The Treaty of Zuhab". Please write a new paragraph with fresh information.} \\
            }
        }
    }
    & \\ \\

    \theutterance \stepcounter{utterance}  
    & \multicolumn{4}{p{0.6\linewidth}}{
        \cellcolor[rgb]{0.9,0.9,0.9}{
            \makecell[{{p{\linewidth}}}]{
                \texttt{\tiny{[P1$\rangle$GM]}}
                \texttt{PARAGRAPH: This military confrontation demonstrated the superior mobility of the eastern cavalry against the heavy infantry of the north, proving that the new dynasty could effectively defend its heartland against a much larger opposing army.} \\
            }
        }
    }
    & & \\ \\

    \theutterance \stepcounter{utterance}  
    & & \multicolumn{4}{p{0.6\linewidth}}{
        \cellcolor[rgb]{0.9,0.9,0.9}{
            \makecell[{{p{\linewidth}}}]{
                \texttt{\tiny{[GM$\rangle$P2]}}
                \texttt{{[}Sentence 9{]}} \\
\texttt{This military confrontation demonstrated the superior mobility of the eastern cavalry against the heavy infantry of the north, proving that the new dynasty could effectively defend its heartland against a much larger opposing army.} \\
            }
        }
    }
    & \\ \\

    \theutterance \stepcounter{utterance}  
    & & & \multicolumn{4}{p{0.6\linewidth}}{
        \cellcolor[rgb]{0.9,0.9,0.9}{
            \makecell[{{p{\linewidth}}}]{
                \texttt{\tiny{[GM$\langle$P2]}}
                \texttt{ANALYSIS: The description highlights a specific tactical advantage: eastern cavalry mobility defeating a larger northern infantry force. This points to a decisive battle where speed and maneuverability were key, likely involving the Mongols or a specific Safavid victory against the Timurids or Ottomans.} \\
\\ 
\texttt{GUESS: The Battle of Chaldiran} \\
            }
        }
    }
     \\ \\

    \theutterance \stepcounter{utterance}  
    & & \multicolumn{4}{p{0.6\linewidth}}{
        \cellcolor[rgb]{0.9,0.9,0.9}{
            \makecell[{{p{\linewidth}}}]{
                \texttt{\tiny{[P1$\langle$GM]}}
                \texttt{Round 9 complete. The Detective's current guess: "The Battle of Chaldiran". Please write a new paragraph with fresh information.} \\
            }
        }
    }
    & \\ \\

    \theutterance \stepcounter{utterance}  
    & \multicolumn{4}{p{0.6\linewidth}}{
        \cellcolor[rgb]{0.9,0.9,0.9}{
            \makecell[{{p{\linewidth}}}]{
                \texttt{\tiny{[P1$\rangle$GM]}}
                \texttt{PARAGRAPH: The decisive engagement at this location resulted in the destruction of the northern army's elite cavalry units, effectively ending their ability to project power into the eastern territories for the remainder of the century.} \\
            }
        }
    }
    & & \\ \\

    \theutterance \stepcounter{utterance}  
    & & \multicolumn{4}{p{0.6\linewidth}}{
        \cellcolor[rgb]{0.9,0.9,0.9}{
            \makecell[{{p{\linewidth}}}]{
                \texttt{\tiny{[GM$\rangle$P2]}}
                \texttt{{[}Sentence 10{]}} \\
\texttt{The decisive engagement at this location resulted in the destruction of the northern army's elite cavalry units, effectively ending their ability to project power into the eastern territories for the remainder of the century.} \\
            }
        }
    }
    & \\ \\

    \theutterance \stepcounter{utterance}  
    & & & \multicolumn{4}{p{0.6\linewidth}}{
        \cellcolor[rgb]{0.9,0.9,0.9}{
            \makecell[{{p{\linewidth}}}]{
                \texttt{\tiny{[GM$\langle$P2]}}
                \texttt{ANALYSIS: The description specifies the destruction of elite northern cavalry units, ending their ability to project power for a century. This confirms the Battle of Chaldiran (1514), where Safavid forces under Ismail I defeated the Timurid army, securing Persian dominance in the region.} \\
\\ 
\texttt{GUESS: The Battle of Chaldiran} \\
            }
        }
    }
     \\ \\

    \theutterance \stepcounter{utterance}  
    & & & \multicolumn{2}{p{0.3\linewidth}}{
        \cellcolor[rgb]{0.9,0.9,0.9}{
            \makecell[{{p{\linewidth}}}]{
                \texttt{\tiny{[GM$|$GM]}}
                \texttt{10} \\
            }
        }
    }
    & & \\ \\

\end{supertabular}
}

\end{document}
